\documentclass[]{jlreq}
\usepackage{bm}
\usepackage{amsmath}
\usepackage[dvipdfmx]{graphicx}
\begin{document}
\title{一端を固定された紐の落下挙動の解析}
\author{黒瀬諭一朗 \thanks{所属：埼玉県立大宮高等学校物理部}}
\maketitle



\section{抄録}
本研究は、一様な重力場の中にある一端を固定された紐の落下挙動を解析したものである。
紐を「$N$個の質点と微小な長さを持つ剛体棒からなる$N$重振り子」として定式化し、
$N\rightarrow\infty$の極限をとることで解析を行った。
その結果、紐は初期状態の直線を保ったまま落下をしようとするが、固定端にデルタ関数的な得意性が現れ、
固定端の直後から紐が急激に折れ曲がりながら落ちていくことがわかった。



\section{研究の背景と目的}
32重振り子のシミュレーション動画を視聴したとき、落下開始時の挙動が紐のように見えた。
この挙動を観察すると、32重振り子は初期状態の直線を保ったまま落下しようとするが、
時間経過と共に固定端側から折れ曲がってゆくことに気がついた。
同時に、紐という連続的な物体がどのように運動するのか疑問に思った。
一端を固定された紐の落下挙動を解明することを目的として、本研究を行った。



\section{方法}
まず、紐を$N$等分することによって、
「$N$個の質点と微小な長さを持つ剛体棒からなる$N$重振り子」として定式化する。
次に、$N$重振り子のラグランジアンを導出し、運動方程式を得る。
最後に、$N\rightarrow\infty$の極限をとり、その結果の物理的意味を考察する。



\section{結果}
\subsection{紐の定式化}



\subsection{ラグランジアンと運動方程式の導出}



\subsection{$N\rightarrow\infty$の極限}



\section{考察}


\begin{thebibliography}{9}
  \item Shiki『N重振り子の運動方程式』
  \item ランダウ=リフシッツ『力学』東京図書
  \item 小出昭一郎『解析力学』岩波書店
\end{thebibliography}
\end{document}
